\documentclass[11pt,a4paper]{article}
\usepackage[T1]{fontenc}
\usepackage[utf8]{inputenc}
\usepackage{amsmath,amssymb,amsthm}
\usepackage[margin=1in]{geometry}
\usepackage[colorlinks=true,linkcolor=blue,urlcolor=blue]{hyperref}
\theoremstyle{plain}
\newtheorem{theorem}{Theorem}
\newtheorem{lemma}[theorem]{Lemma}
\newtheorem{proposition}[theorem]{Proposition}
\newtheorem{corollary}[theorem]{Corollary}
\theoremstyle{definition}
\newtheorem{remark}[theorem]{Remark}
% A quoted conjecture can be a single hundred-term formula whose terms are thirty-digit
% numbers. TeX cannot hyphenate those, so without a little stretch every such quote runs off
% the page; the linked-a-file papers are all of that shape.
\setlength{\emergencystretch}{3em}

\title{The conjectured P-recursive recurrence for OEIS A292461, proved}
\author{Adrian Perez Fontelles\\ \small Independent researcher}
\date{15 September 2026}
\begin{document}
\maketitle

\begin{abstract}
OEIS A292461 carries a conjectured recurrence of order $4$ whose coefficients are polynomials in
$n$ rather than constants. It is true, and it is decidable rather than empirical. The entry
states its generating function as a fact, and that function is algebraic; the conjecture then
becomes the vanishing of an explicit element of an algebraic function field, which is settled
symbolically. The result is an identity, not a run of verified terms: it holds at every index
at once.
\end{abstract}

\noindent\small 2020 Mathematics Subject Classification. 05A15, 33F10, 11B37.\normalsize

\section{The conjecture}

OEIS A292461 is ``Expansion of (1 - x - x\^{}2 + sqrt((1 - x - x\^{}2)\^{}2 - 4*x\^{}3))/2 in powers of x.''. It has offset $0$ and begins
\[
1, -1, -1, -1, -1, -2, -4, -8,\ \dots
\]
The entry states, as a conjecture contributed by its author and never marked settled:
\begin{quote}\raggedright\small
Conjecture D-finite with recurrence: n*a(n) +(-2*n+3)*a(n-1) +(-n+3)*a(n-2) +(-2*n+9)*a(n-3) +(n-6)*a(n-4)=0. - \_R. J. Mathar\_, Jan 24 2020
\end{quote}
As of the ``Last modified'' line on the live entry (Sep 08 2022 08:46:19, revision 26)
this is still recorded as a conjecture, and nothing on the entry records it as proved.

\section{The generating function}

The entry gives, as a statement of fact and not as a conjecture,
\[
A(x)\;=\;\frac{- 2 x^{3} + \left(- x^{2} - x + 1\right) \left(- x^{2} - x + \sqrt{\left(x^{2} - 3 x + 1\right) \left(x^{2} + x + 1\right)} + 1\right)}{- x^{2} - x + \sqrt{\left(x^{2} - 3 x + 1\right) \left(x^{2} + x + 1\right)} + 1} .
\]
This is an algebraic function. Before it is used for anything it is checked against the entry's
own data: its Taylor coefficients reproduce every one of the $33$ published terms
exactly, in exact arithmetic. That check is not a formality --- on A116388 the generating
function an entry states as fact does NOT generate the terms it publishes, and a proof built on
it would have been a proof from a false premise.

\section{A P-recursive claim is the vanishing of one function}

\begin{lemma}\label{lem:theta}
Let $\theta=x\,\frac{d}{dx}$, so $\theta$ acts on $x^m$ as multiplication by $m$. For a
polynomial $p$ and an integer $i\ge0$,
\[
\sum_{n} p(n)\,a(n-i)\,x^{n}\;=\;x^{i}\,\bigl(p(\theta+i)A\bigr)(x).
\]
\end{lemma}

\begin{proof}
Write $m=n-i$. Then $\sum_m p(m+i)a(m)x^{m+i}=x^i\sum_m p(m+i)a(m)x^m$, and $p(\theta+i)$
applied to $A=\sum_m a(m)x^m$ multiplies the $m$-th coefficient by $p(m+i)$.
\end{proof}

Put $b(n)=\sum_{i=0}^{4}p_i(n)\,a(n-i)$, the left-hand side of the conjecture, and
$B(x)=\sum_n b(n)x^n$. By Lemma~\ref{lem:theta},
\[
B(x)\;=\;\sum_{i=0}^{4}x^{i}\,\bigl(p_i(\theta+i)A\bigr)(x),
\]
an explicit element of the algebraic function field generated by $A$. The conjecture asserts
$b(n)=0$; so it holds for all $n$ exactly when $B=0$, and for all $n>d$ exactly when $B$ is a
polynomial of degree $d$. Both are decided symbolically, with no truncation and no numerics.

\section{The theorem}

\begin{theorem}
\[
\left(n\right)a(n) + \left(3 - 2 n\right)a(n-1) + \left(3 - n\right)a(n-2) + \left(9 - 2 n\right)a(n-3) + \left(n - 6\right)a(n-4)\;=\;0
\]
for every $n>5$.
\end{theorem}

\begin{proof}
$B$ was formed as above and simplified in the function field. Here $B$ is a polynomial of degree $5$. A polynomial contributes only to the coefficients of $x^k$ with $k\le5$, so $b(n)=0$ for every $n>5$. Since $A$ reproduces the
entry's published terms, $b$ is the sequence the entry's conjecture is about, and the statement
follows.
\end{proof}

\section{Verification}

First, the generating function reproduces all $33$ published terms exactly. Second,
the conjectured recurrence was evaluated directly on those terms, with no generating function
involved, and holds wherever the proved range and the data overlap. Third, the residual $B$ was
computed symbolically rather than from a truncated series, and its series expansion was taken
independently to confirm the coefficients it predicts.

\begin{thebibliography}{9}
\bibitem{oeis} The OEIS Foundation, \emph{The On-Line Encyclopedia of Integer Sequences},
\url{https://oeis.org}, 2026. Sequence A292461.
\bibitem{stanley2} R.~P.~Stanley, \emph{Enumerative Combinatorics}, Volume 2, Cambridge
University Press, 1999. (Algebraic and D-finite generating functions.)
\bibitem{kp} M.~Kauers and P.~Paule, \emph{The Concrete Tetrahedron}, Springer, 2011.
\end{thebibliography}

\end{document}
